O algoritmo Shell Sort apresenta um comportamento diferente dos demais algoritmos analisados, pois sua complexidade depende da sequência de \textit{gaps} utilizada. Neste trabalho, o \textit{gap} inicial é obtido pela divisão do tamanho do vetor pelo divisor escolhido e, a cada etapa, é novamente dividido pelo mesmo divisor até atingir o valor 1.

Na implementação utilizada nos experimentos, foi adotado o divisor igual a 2. Dessa forma, para um vetor de tamanho \(n\), a sequência de \textit{gaps} é obtida aproximadamente por:

\begin{equation}
\frac{n}{2},\frac{n}{4},\frac{n}{8},\ldots,1
\end{equation}

Por exemplo, para um vetor de tamanho 64, utilizando o divisor 2, a sequência de \textit{gaps} será:

\begin{equation}
32,16,8,4,2,1
\end{equation}

Já utilizando o divisor 4, a sequência será:

\begin{equation}
16,4,1
\end{equation}

Assim, um divisor maior tende a reduzir a quantidade de etapas realizadas. Entretanto, isso não significa necessariamente que o algoritmo será mais rápido, pois o desempenho também depende da quantidade de comparações e deslocamentos realizados em cada etapa e da desordem restante quando o \textit{gap} atinge o valor 1.

\subsubsection{Melhor caso}

O melhor caso ocorre quando o vetor já está ordenado de forma crescente. Nessa situação, os elementos já estão nas posições adequadas, de modo que as etapas de ordenação não precisam realizar deslocamentos significativos.

Em cada etapa, o algoritmo percorre os elementos do vetor, realizando comparações para verificar se a ordem está correta. Considerando uma sequência de \textit{gaps} obtida pela divisão sucessiva por uma constante maior que 1, o número de etapas é proporcional a \(\log n\).

Como cada etapa realiza aproximadamente \(n\) verificações no vetor, o custo pode ser estimado por:

\begin{equation}
T(n) = O(n\log n)
\end{equation}

Essa expressão representa uma estimativa para o número de operações nesse cenário, considerando a sequência de \textit{gaps} adotada. A complexidade exata deve levar em conta os detalhes da implementação e a sequência utilizada.

\subsubsection{Pior caso}

No pior caso, a disposição inicial dos elementos pode exigir uma grande quantidade de comparações e deslocamentos. Embora as etapas com \textit{gaps} maiores reduzam parte da desordem, a etapa final, com \textit{gap} igual a 1, pode exigir um número elevado de operações.

Nessa etapa, o Shell Sort realiza uma ordenação equivalente ao Insertion Sort sobre o vetor inteiro. No pior caso, a quantidade de operações dessa etapa pode ser representada por:

\begin{equation}
T(n) = 1+2+3+\cdots+(n-1)
\end{equation}

Aplicando a fórmula da soma dos termos de uma progressão aritmética, temos:

\begin{equation}
T(n) = \frac{n(n-1)}{2}
\end{equation}

Desenvolvendo a expressão:

\begin{equation}
T(n) = \frac{n^2-n}{2}
\end{equation}

O termo de maior crescimento é \(n^2\). Portanto, a etapa final pode apresentar custo quadrático:

\begin{equation}
T(n) = \Theta(n^2)
\end{equation}

Entretanto, essa análise considera o custo máximo da etapa com \textit{gap} igual a 1. Para determinar a complexidade de pior caso do Shell Sort completo, é necessário considerar também o efeito das etapas anteriores e a sequência de \textit{gaps} utilizada.

\subsubsection{Caso médio}

O caso médio considera entradas com diferentes disposições dos elementos, como vetores gerados aleatoriamente. Diferentemente de algoritmos como Insertion Sort, Bubble Sort e Selection Sort, cuja análise pode ser apresentada sem depender de uma sequência de intervalos, o comportamento do Shell Sort depende diretamente da sequência de \textit{gaps} adotada.

Alterar o divisor modifica a quantidade de etapas e os intervalos utilizados nas comparações. Por exemplo, para um vetor de tamanho 64, o divisor 2 produz a sequência:

\begin{equation}
32,16,8,4,2,1
\end{equation}

Enquanto o divisor 4 produz:

\begin{equation}
16,4,1
\end{equation}

Apesar de a segunda sequência possuir menos etapas, isso não garante um menor tempo de execução. O desempenho depende de como cada etapa reorganiza os elementos e de quanto trabalho permanece para a etapa final.

Dessa forma, não é adequado atribuir uma única complexidade média ao Shell Sort sem especificar a sequência de \textit{gaps} e as condições consideradas na análise. Neste trabalho, o caso médio será avaliado experimentalmente por meio dos tempos de execução obtidos para vetores com elementos em ordem aleatória.

\subsubsection{Resumo das complexidades}

\begin{table}[H]
    \centering
    \caption{Análise da complexidade do algoritmo Shell Sort}
    \label{tab:complexidade-shell}
    \begin{tabular}{lc}
        \hline
        \textbf{Caso} & \textbf{Complexidade de tempo} \\
        \hline
        Melhor caso & Depende da sequência de \textit{gaps} \\
        Caso médio & Depende da sequência de \textit{gaps} \\
        Pior caso & Depende da sequência de \textit{gaps} \\
        \hline
    \end{tabular}
\end{table}

Conclui-se que o Shell Sort apresenta uma complexidade que depende da sequência de \textit{gaps} utilizada. O divisor determina como os intervalos são reduzidos e influencia a quantidade de etapas e de operações realizadas durante a ordenação.

Neste trabalho, foi utilizado o divisor 2 nos experimentos, produzindo uma sequência de \textit{gaps} obtida pela divisão sucessiva por 2 até atingir o valor 1. Uma possível extensão dos experimentos seria comparar diferentes divisores, permitindo investigar como a escolha da sequência de \textit{gaps} influencia o tempo de execução do algoritmo.